%%%%%%%%%%%%%%%%%%%%%%% file template.tex %%%%%%%%%%%%%%%%%%%%%%%%%
%
% This is a general template file for the LaTeX package SVJour3
% for Springer journals.          Springer Heidelberg 2010/09/16
%
% Copy it to a new file with a new name and use it as the basis
% for your article. Delete % signs as needed.
%
% This template includes a few options for different layouts and
% content for various journals. Please consult a previous issue of
% your journal as needed.
%
%%%%%%%%%%%%%%%%%%%%%%%%%%%%%%%%%%%%%%%%%%%%%%%%%%%%%%%%%%%%%%%%%%%
%
% First comes an example EPS file -- just ignore it and
% proceed on the \documentclass line
% your LaTeX will extract the file if required
\begin{filecontents*}{example.eps}
%!PS-Adobe-3.0 EPSF-3.0
%%BoundingBox: 19 19 221 221
%%CreationDate: Mon Sep 29 1997
%%Creator: programmed by hand (JK)
%%EndComments
gsave
newpath
  20 20 moveto
  20 220 lineto
  220 220 lineto
  220 20 lineto
closepath
2 setlinewidth
gsave
  .4 setgray fill
grestore
stroke
grestore
\end{filecontents*}
%
\RequirePackage{fix-cm}
%
\documentclass[smallcondensed, draft, natbib]{svjour3}
\smartqed  % flush right qed marks, e.g. at end of proof
%
\usepackage{graphicx}
%
% \usepackage{mathptmx}      % use Times fonts if available on your TeX system
%
% insert here the call for the packages your document requires
%\usepackage{latexsym}
% etc.
%
% please place your own definitions here and don't use \def but
% \newcommand{}{}
%
% Insert the name of "your journal" with
\journalname{myjournal}
%
%
% % header.tex
% % this is where you load pacakges, specify custom formats, etc.
%
% \usepackage[left=1in,right=1in,top=0.75in,footskip=25pt]{geometry}
% % \usepackage{changepage}
% \usepackage{amsmath,amsthm,amssymb,amsfonts}
% \usepackage{mathtools}
% \usepackage{mathrsfs}
% \usepackage{bbm}
% % enumitem for custom lists
% \usepackage{enumitem}
% % Load dsfont this to get proper indicator function (bold 1) with \mathds{1}:
% \usepackage{dsfont}
% \usepackage{centernot}
%
\usepackage[usenames,dvipsnames]{xcolor}
% \usepackage[parfill]{parskip}
%
% % set up graphics
% \usepackage{graphicx}
%
% \usepackage[sorting=nyt,backend=biber,bibstyle=alphabetic,citestyle=alphabetic,giveninits=true]{biblatex}
%
% \usepackage{fancyhdr}
% \pagestyle{fancy}
% \setlength{\headheight}{40pt}
%
% %%%%%%%%%%%%%%%%%%%%%%%%%%%%%%%%%%%%%%%%%%%%%%%%%%%%%%%%%%%%%%%%%%%%%%%%%%%%%%%%%%%%
% % most other packages you might use should be loaded before hyperref
% %%%%%%%%%%%%%%%%%%%%%%%%%%%%%%%%%%%%%%%%%%%%%%%%%%%%%%%%%%%%%%%%%%%%%%%%%%%%%%%%%%%%
%
% % Set up hyperlinks:
% \definecolor{RefColor}{rgb}{0,0,.65}
% \usepackage[colorlinks,linkcolor=RefColor,citecolor=RefColor,urlcolor=RefColor]{hyperref}
%
% \usepackage[capitalize]{cleveref}
% \crefname{appsec}{Appendix}{Appendices} % you can tell cleveref what to call things

% % bf series
% \def\bfA{\mathbf{A}}
% \def\bfB{\mathbf{B}}
% \def\bfC{\mathbf{C}}
% \def\bfD{\mathbf{D}}
% \def\bfE{\mathbf{E}}
% \def\bfF{\mathbf{F}}
% \def\bfG{\mathbf{G}}
% \def\bfH{\mathbf{H}}
% \def\bfI{\mathbf{I}}
% \def\bfJ{\mathbf{J}}
% \def\bfK{\mathbf{K}}
% \def\bfL{\mathbf{L}}
% \def\bfM{\mathbf{M}}
% \def\bfN{\mathbf{N}}
% \def\bfO{\mathbf{O}}
% \def\bfP{\mathbf{P}}
% \def\bfQ{\mathbf{Q}}
% \def\bfR{\mathbf{R}}
% \def\bfS{\mathbf{S}}
% \def\bfT{\mathbf{T}}
% \def\bfU{\mathbf{U}}
% \def\bfV{\mathbf{V}}
% \def\bfW{\mathbf{W}}
% \def\bfX{\mathbf{X}}
% \def\bfY{\mathbf{Y}}
% \def\bfZ{\mathbf{Z}}
% \def\bfLambda{\mathbf{\Lambda}}
% \def\bfXi{\mathbf{\Xi}}
%
% \def\bfa{\mathbf{a}}
% \def\bfb{\mathbf{b}}
% \def\bfc{\mathbf{c}}
% \def\bfd{\mathbf{d}}
% \def\bfe{\mathbf{e}}
% \def\bff{\mathbf{f}}
% \def\bfg{\mathbf{g}}
% \def\bfh{\mathbf{h}}
% \def\bfi{\mathbf{i}}
% \def\bfj{\mathbf{j}}
% \def\bfk{\mathbf{k}}
% \def\bfl{\mathbf{l}}
% \def\bfm{\mathbf{m}}
% \def\bfn{\mathbf{n}}
% \def\bfo{\mathbf{o}}
% \def\bfp{\mathbf{p}}
% \def\bfq{\mathbf{q}}
% \def\bfr{\mathbf{r}}
% \def\bfs{\mathbf{s}}
% \def\bft{\mathbf{t}}
% \def\bfu{\mathbf{u}}
% \def\bfv{\mathbf{v}}
% \def\bfw{\mathbf{w}}
% \def\bfx{\mathbf{x}}
% \def\bfy{\mathbf{y}}
% \def\bfz{\mathbf{z}}
%
% \def\bfzero{\mathbf{0}}
% \def\bfone{\mathbf{1}}
%
% % bb series
% \def\bbA{\mathbb{A}}
% \def\bbB{\mathbb{B}}
% \def\bbC{\mathbb{C}}
% \def\bbD{\mathbb{D}}
% \def\bbE{\mathbb{E}}
% \def\bbF{\mathbb{F}}
% \def\bbG{\mathbb{G}}
% \def\bbH{\mathbb{H}}
% \def\bbI{\mathbb{I}}
% \def\bbJ{\mathbb{J}}
% \def\bbK{\mathbb{K}}
% \def\bbL{\mathbb{L}}
% \def\bbM{\mathbb{M}}
% \def\bbN{\mathbb{N}}
% \def\bbO{\mathbb{O}}
% \def\bbP{\mathbb{P}}
% \def\bbQ{\mathbb{Q}}
% \def\bbR{\mathbb{R}}
% \def\bbS{\mathbb{S}}
% \def\bbT{\mathbb{T}}
% \def\bbU{\mathbb{U}}
% \def\bbV{\mathbb{V}}
% \def\bbW{\mathbb{W}}
% \def\bbX{\mathbb{X}}
% \def\bbY{\mathbb{Y}}
% \def\bbZ{\mathbb{Z}}
%
% % cal series
% \def\calA{\mathcal{A}}
% \def\calB{\mathcal{B}}
% \def\calC{\mathcal{C}}
% \def\calD{\mathcal{D}}
% \def\calE{\mathcal{E}}
% \def\calF{\mathcal{F}}
% \def\calG{\mathcal{G}}
% \def\calH{\mathcal{H}}
% \def\calI{\mathcal{I}}
% \def\calJ{\mathcal{J}}
% \def\calK{\mathcal{K}}
% \def\calL{\mathcal{L}}
% \def\calM{\mathcal{M}}
% \def\calN{\mathcal{N}}
% \def\calO{\mathcal{O}}
% \def\calP{\mathcal{P}}
% \def\calQ{\mathcal{Q}}
% \def\calR{\mathcal{R}}
% \def\calS{\mathcal{S}}
% \def\calT{\mathcal{T}}
% \def\calU{\mathcal{U}}
% \def\calV{\mathcal{V}}
% \def\calW{\mathcal{W}}
% \def\calX{\mathcal{X}}
% \def\calY{\mathcal{Y}}
% \def\calZ{\mathcal{Z}}
%
% % scr series
% \def\scrA{\mathscr{A}}
% \def\scrB{\mathscr{B}}
% \def\scrC{\mathscr{C}}
% \def\scrD{\mathscr{D}}
% \def\scrE{\mathscr{E}}
% \def\scrF{\mathscr{F}}
% \def\scrG{\mathscr{G}}
% \def\scrH{\mathscr{H}}
% \def\scrI{\mathscr{I}}
% \def\scrJ{\mathscr{J}}
% \def\scrK{\mathscr{K}}
% \def\scrL{\mathscr{L}}
% \def\scrM{\mathscr{M}}
% \def\scrN{\mathscr{N}}
% \def\scrO{\mathscr{O}}
% \def\scrP{\mathscr{P}}
% \def\scrQ{\mathscr{Q}}
% \def\scrR{\mathscr{R}}
% \def\scrS{\mathscr{S}}
% \def\scrT{\mathscr{T}}
% \def\scrU{\mathscr{U}}
% \def\scrV{\mathscr{V}}
% \def\scrW{\mathscr{W}}
% \def\scrX{\mathscr{X}}
% \def\scrY{\mathscr{Y}}
% \def\scrZ{\mathscr{Z}}
%
% \newtheorem{theorem}{Theorem}
% \newtheorem{lemma}{Lemma}
% \newtheorem{proposition}{Proposition}
% \newtheorem{corollary}{Corollary}
% \newtheorem{observation}{Observation}
% \newtheorem{definition}{Definition}
% \newtheorem{claim}{Claim}
% \newtheorem{fact}{Fact}
% \newtheorem{assumption}{Assumption}
% \newtheorem{example}{Example}
%
%
% %%%%%%%%%%%%%%%%%%%%%%%%%%%%%%%%%%%%%%%%%%%%%%%%%%%%%%%%%%
% % text short-cuts
% \def\iid{i.i.d.\ } %i.i.d.
% \def\ie{i.e.\ }
% \def\eg{e.g.\ }
%
% % operators
% \def\P{\bbP} %fundamental probability
% \def\E{\bbE} %expectation
% \newcommand{\Cov}[1]{\text{Cov}\left(#1\right)}
% \newcommand{\Var}[1]{\text{Var}\left(#1\right)}
%
% % conditional expectation
% \DeclarePairedDelimiterX\bigCond[2]{[}{]}{#1 \;\delimsize\vert\; #2}
% \newcommand{\conditional}[3][]{\bbE_{#1}\bigCond*{#2}{#3}}
% \def\indicator{\mathbbm{1}} % indicator function
%
% \DeclareMathOperator*{\argmax}{arg\,max}
% \DeclareMathOperator*{\argmin}{arg\,min}
%
% colored texts
\newcommand{\red}[1]{\textcolor{red}{#1}}
\newcommand{\blue}[1]{\textcolor{blue}{#1}}
%
% \newcommand{\expectationof}[1]{\bbE{\left[#1\right]}}
% \newcommand{\defeq}{\vcentcolon=}
%
%
% \newcommand{\x}{x}
% \newcommand{\xk}{x_{k}}
% \newcommand{\mxk}{\bar{x}_{k}}
% \newcommand{\mxt}{\bar{x}_{T}}
% \newcommand{\xkk}{x_{k+1}}
% \newcommand{\xopt}{x^{*}}
% \newcommand{\etak}{\eta_{k}}
%
% \newcommand{\norm}[1]{\left\|#1\right\|}
% \newcommand{\normsq}[1]{\left\|#1\right\|^{2}}
% \newcommand{\tr}[1]{\,\mathrm{tr}\,#1}
% \newcommand{\eye}{\bfI}
% \newcommand{\inner}[2]{\left\langle#1\,,\,#2\right\rangle}
% \newcommand{\dom}{\text{dom}}
%
% \newcommand{\fof}[1]{f(#1)}
% \newcommand{\gradf}[1]{\nabla f(#1)}
% \newcommand{\sgradf}[2]{\nabla f_{#1}(#2)}
% \newcommand{\hessf}[1]{\nabla^2 f(#1)}
% \newcommand{\invhessf}[1]{\nabla^2 f(#1)^{-1}}
% \newcommand{\shessf}[2]{\mathbf{H}_{#1}(#2)}
% \newcommand{\spanof}[1]{\text{span}\left\{#1\right\}}
%
% \DeclarePairedDelimiter\ceil{\lceil}{\rceil}
% \DeclarePairedDelimiter\floor{\lfloor}{\rfloor}


\begin{document}

\title{Newton-Laplace Updates for Block Coordinate Descent\thanks{\red{Grants or other notes
about the article that should go on the front page should be
placed here. General acknowledgments should be placed at the end of the article.}}
}

%\titlerunning{Short form of title}        % if too long for running head

\author{Si Yi Meng         \and
        Mark Schmidt %etc.
}

%\authorrunning{Short form of author list} % if too long for running head

\institute{Si Yi Meng \at
              Department of Computer Science, The University of British Columbia \\
              \email{mengxixi@cs.ubc.ca}           %  \\
%             \emph{Present address:} of F. Author  %  if needed
           \and
           Mark Schmidt \at
           Department of Computer Science, The University of British Columbia \\
           \email{schmidtm@cs.ubc.ca}
}

\date{Received: date / Accepted: date}
% The correct dates will be entered by the editor


\maketitle

\begin{abstract}
\red{Insert your abstract here. Include keywords, PACS and mathematical
subject classification numbers as needed.}
\keywords{\red{First keyword \and Second keyword \and More}}
% \PACS{PACS code1 \and PACS code2 \and more}
% \subclass{MSC code1 \and MSC code2 \and more}
\end{abstract}

\section{Introduction}
\label{intro}
Your text comes here. Separate text sections with
\section{Section title}
\label{sec:1}
Text with citations %\cite{RefB} and \cite{RefJ}.
\subsection{Subsection title}
\label{sec:2}
as required. Don't forget to give each section
and subsection a unique label (see Sect.~\ref{sec:1}).
\paragraph{Paragraph headings} Use paragraph headings as needed.
\begin{equation}
a^2+b^2=c^2
\end{equation}

% For one-column wide figures use
\begin{figure}
% Use the relevant command to insert your figure file.
% For example, with the graphicx package use
  \includegraphics{example.eps}
% figure caption is below the figure
\caption{Please write your figure caption here}
\label{fig:1}       % Give a unique label
\end{figure}
%
% For two-column wide figures use
\begin{figure*}
% Use the relevant command to insert your figure file.
% For example, with the graphicx package use
  \includegraphics[width=0.75\textwidth]{example.eps}
% figure caption is below the figure
\caption{Please write your figure caption here}
\label{fig:2}       % Give a unique label
\end{figure*}
%
% For tables use
\begin{table}
% table caption is above the table
\caption{Please write your table caption here}
\label{tab:1}       % Give a unique label
% For LaTeX tables use
\begin{tabular}{lll}
\hline\noalign{\smallskip}
first & second & third  \\
\noalign{\smallskip}\hline\noalign{\smallskip}
number & number & number \\
number & number & number \\
\noalign{\smallskip}\hline
\end{tabular}
\end{table}


\begin{acknowledgements}
\red{If you'd like to thank anyone, place your comments here and remove the percent signs.}
\end{acknowledgements}

% Authors must disclose all relationships or interests that
% could have direct or potential influence or impart bias on
% the work:
%
\section*{Conflict of interest}
%
\red{The authors declare that they have no conflict of interest.}


% BibTeX users please use one of
\bibliographystyle{spbasic}      % basic style, author-year citations
% \bibliographystyle{spmpsci}      % mathematics and physical sciences
%\bibliographystyle{spphys}       % APS-like style for physics
\bibliography{refs}   % name your BibTeX data base


\end{document}
% end of file template.tex
